\documentclass[10pt,a4paper]{amsart}
\usepackage[margin=19mm]{geometry}
\usepackage{amsmath,amssymb,mathtools}
\usepackage[scr=rsfs,cal=euler]{mathalfa}
\usepackage{xfrac}
\usepackage[unicode,hidelinks,bookmarks=false]{hyperref}
\newcommand{\ie}{i.e.\ }
\newcommand{\cf}{cf.\ }
\newcommand{\resp}{respectively\ }
\newcommand{\Subsection}{\subsection}
\providecommand{\nobreakdash}{\nobreak\hbox{-}}
\setlength{\emergencystretch}{2em}
\multlinegap0pt
\usepackage{algorithm}\usepackage{algpseudocode}
\usepackage[shortlabels]{enumitem}
\usepackage{siunitx}
\usepackage{tikz}
\usepackage{dsfont}
\usepackage{mathtools}
\newtheorem{theorem}{Theorem}
\newtheorem{pf}{Pf}
\title{Training-free Adjustable Polynomial Graph Filtering for Ultra-fast Multimodal Recommendation}
\author{Yu-Seung Roh, Joo-Young Kim, Jin-Duk Park, Won-Yong Shin}
\date{Source: arXiv:2503.04406v3 (CC BY 4.0)}
\begin{document}
\maketitle

\noindent\emph{Source and licence.} Training-free Adjustable Polynomial Graph Filtering for Ultra-fast Multimodal Recommendation. Version arXiv:2503.04406v3 (CC BY 4.0), distributed by arXiv under the Creative Commons Attribution 4.0 International licence, http://creativecommons.org/licenses/by/4.0/, as recorded in the arXiv licence metadata for this exact version and shown on the abstract page https://arxiv.org/abs/2503.04406v3. Full licence notice: \texttt{licenses/arxiv-2503.04406-CC-BY-4.0.txt}.\par\smallskip

\noindent\textit{Editorial scope.} Counted unit: the article's theorem thm\_poly (source0.tex lines 564--573). The article's complete original proof of it is delivered with it (source0.tex lines 575--597, one \texttt{proof} environment of 23 lines). No mathematics is written, repaired or re-derived: the statement and the proof are the article's own lines, in the article's own notation, and no retained line is substituted or re-flowed.  No further source-local context is needed: every cross-reference of every retained line resolves inside the two delivered spans.  Nothing was omitted from any retained span, so the package carries no omission record.  No retained line contains a citation, so the wrapper carries no bibliography and no bibliography entry.  No retained line contains a figure environment, an image-inclusion command or drawing code, so no image is delivered and the package declares no asset dependency.  Editorial formatting only, in addition: an AMS \texttt{amsart} preamble that restates the article's own declarations and macros used by the retained lines, namely the environments \texttt{theorem}, \texttt{pf} on separate counters, together with the standard packages those lines need.  The retained environments print in this document as Theorem 1, Pf 1 in order of appearance, whereas the article prints the same items with the numbering of its own full document.  The complete native source of the article as distributed in the arXiv e-print for this exact version is bundled unchanged as \texttt{sources/174325/source0.tex}.
\begin{theorem}
    \label{thm_poly}
     For any admissible range of eigenvalues, the matrix polynomial can be expressed as:

     \begin{equation}
     \label{poly_filter}
         \bar P ^f_\star = \sum_{k=1}^{K} \frac{a_k}{(\lambda^\star)^{k-1}} (\bar{P}_\star-\lambda_{min}^\star  I)^k,
     \end{equation}
     where $\bar P^f_\star$ is the graph filter associated with the similarity graph $\bar{P}_\star$, having the frequency response function of $h(\lambda)=\sum_{k=1}^K\frac{a_k}{(\lambda^\star)^{k-1}} (\lambda^\star-\lambda)^k$, and $\lambda^\star = \lambda^\star_{max} - \lambda^\star_{min}$. Here, $\lambda^\star_{min}$ and $\lambda^\star_{min}$ are the minimum and maximum eigenvalues of $\bar P_\star$, respectively.
\end{theorem}

\begin{pf}
    First, we define the graph Laplacian $L_\star$ for a specific modality $\star$ as follows:

    \begin{equation}
        \bar L_\star = \lambda_{max}^\star I - \bar P_\star = U\Lambda^\star U^T,
    \end{equation}
    where $\Lambda^\star$ is a set of eigenvalues of $L_\star$. Since the eigenvalues of $\bar P_\star$ fall within $[\lambda_{min}^\star, \lambda_{max}^\star]$, those of $L_\star$ lie within $[0, \lambda^\star]$. Based on the graph Laplacian $L_\star$, we have

    \begin{equation}
        \begin{aligned}
            \bar P_\star - \lambda^\star_{min}I &= \lambda^\star_{max}I-L_\star - \lambda^\star_{min}I \\ &= U(\lambda^\star_{max}I - \Lambda^\star - \lambda^\star_{min}I)U^T \\ &= U(\lambda^\star I -\Lambda^\star)U^T,
        \end{aligned}
    \end{equation}
    which implies that $\bar P_\star - \lambda^\star_{min}I$ is a graph filter having a frequency response function of $h(\lambda)=\lambda^\star-\lambda$ for $[0, \lambda^\star]$. This corresponds to a linear LPF whose frequency response defined over the interval $[0,\lambda^\star]$. Now, we consider the $k$th-order polynomial filter. Then, we have

    \begin{equation}
        \begin{aligned}
            \frac{1}{(\lambda^\star)^{k-1}}(\bar P_\star-\lambda^\star_{min}I)^k &=\frac{1}{(\lambda^\star)^{k-1}} \underbrace {U (\lambda^\star I-\Lambda^\star) U^T \dots U(\lambda^\star I-\Lambda^\star) U^T}_{k\ \text{times}} \\&=U\left( \frac{(\lambda^\star I -\Lambda^\star)^k}{(\lambda^\star)^{k-1}} \right) U^T,
        \end{aligned}
    \end{equation}
    which indicates a graph filter with a frequency response function of $h(\lambda) = \frac{(\lambda^\star - \lambda)^k}{(\lambda^\star)^{k-1}}$ for $\lambda\in[0, \lambda^\star]$. Since $h(0) = \lambda^\star$, $h(\lambda^\star)=0$, and $h(\lambda)$ is a monotonically decreasing function in the interval $[0, \lambda^\star]$, the filtered value is still bounded in $[0,\lambda^\star]$. Therefore, the proposed filter design guarantees spectral boundedness of $\bar{P}_\star^f$ within $[0, \lambda^\star]$, even when eigenvalues of $\bar{P}_\star$ exceed the standard unit interval. This completes the proof of this theorem.
    \hfill $\square$
\end{pf}

\end{document}
